\documentclass[11pt,a4paper]{article}

\input{../../sablony/pripravy/priprava.sty}

\nastavPripravu{4. ročník}{Kombinatorika}{38}{Permutace s opakováním}

\begin{document}

\blok{Motivace a návaznost}

Kolik různých slov (i nesmyslných) vznikne přeskupením písmen slova MAMA? Stejná písmena nelze vzájemně rozlišit.

\poznamka{Předpoklady: permutace různých prvků, faktoriály a kombinace. Cíl: vysvětlit pevné četnosti stejných prvků a zdůvodnit dělení faktoriály četností.}

\blok{Výklad a odvození}

\begin{definiceBox}
\textbf{Permutace s opakováním} uspořádávají všech $n$ prvků, které mohou zahrnovat nerozlišitelné kopie stejných druhů. Mají-li jednotlivé druhy pevné četnosti $n_1,\ldots,n_r$, kde $n_1+\cdots+n_r=n$, potom
\[
P(n;n_1,\ldots,n_r)=\frac{n!}{n_1!\cdots n_r!}.
\]
\end{definiceBox}
Nejprve dočasně očíslujeme stejné kopie a seřadíme všech $n$ odlišných objektů. Každé viditelné pořadí se objeví přesně $n_1!\cdots n_r!$-krát. Zde \textbf{nevolíme libovolnou četnost}; počty jednotlivých symbolů jsou předepsány.

\blok{Společně řešené příklady}

\begin{prikladReseniBox}[title={Příklad 1}]
Kolik různých přeskupení písmen má slovo MAMA?
\tcblower
Dvě M a dvě A: $\frac{4!}{2!2!}=6$.
\end{prikladReseniBox}



\begin{prikladReseniBox}[title={Příklad 2}]
Kolik různých přeskupení má slovo BANANA?
\tcblower
Šest písmen: A třikrát, N dvakrát, B jednou. $\frac{6!}{3!2!}=60$.
\end{prikladReseniBox}



\begin{prikladReseniBox}[title={Příklad 3}]
Kolik řad vznikne ze tří červených a dvou modrých nerozlišitelných korálků?
\tcblower
$\frac{5!}{3!2!}=10$.
\end{prikladReseniBox}



\blok{Samostatná práce}

\begin{zadaniReseniBox}[title={Samostatná práce 1}]
Urči počet různých přeskupení slova TATA.
\tcblower
$\frac{4!}{2!2!}=6$.
\end{zadaniReseniBox}



\begin{zadaniReseniBox}[title={Samostatná práce 2}]
Kolik pořadí vznikne ze čtyř bílých a dvou černých nerozlišitelných karet?
\tcblower
$\frac{6!}{4!2!}=15$.
\end{zadaniReseniBox}



\begin{zadaniReseniBox}[title={Samostatná práce 3}]
Kolik různých řetězců délky 5 obsahuje přesně tři nuly a dvě jedničky?
\tcblower
$\frac{5!}{3!2!}=10$.
\end{zadaniReseniBox}



\blok{Rozšíření}

\begin{zadaniReseniBox}[title={Rozšíření 1}]
Z číslic $0,0,1,1,2$ vytvoříme pěticiferná čísla, použijeme-li všechny číslice. Kolik jich je?
\tcblower
Všech přeskupení je $5!/(2!2!)=30$. S nulou na začátku je $4!/(1!2!)=12$. Platných pěticiferných čísel je $30-12=18$.
\end{zadaniReseniBox}



\blok{Domácí úkol}

\poznamka{Doporučená práce doma: přibližně 15--30 minut. Úlohy jsou seřazeny od základních k náročnějším; učitel může poslední podúlohu označit jako volitelnou.}

\begin{zadaniReseniBox}[title={Domácí úkol}]

\begin{enumerate}[label=\alph*),leftmargin=*,itemsep=2pt]

\item Kolik různých přeskupení má slovo ANANAS?
\item Kolik různých řad ze tří červených a tří modrých nerozlišitelných kostek?
\item Z číslic $1,1,2,2,3$ sestav všechna různá pěticiferná čísla; urči jejich počet.
\item Kolik různých přeskupení má slovo AABBC?
\item Vysvětli, proč do počtu permutací s opakováním slova AABB nepatří nové slovo, pokud prohodíme obě A.
\item Kolik řad vytvoří tři stejné červené a dvě stejné modré kostky, pokud všechny červené musejí stát pohromadě?

\end{enumerate}

\tcblower

\begin{enumerate}[label=\alph*),leftmargin=*,itemsep=2pt]

\item $6!/(3!2!)=60$.
\item $6!/(3!3!)=20$.
\item $5!/(2!2!)=30$.
\item $5!/(2!2!)=30$.
\item Obě A jsou nerozlišitelné, jejich záměna nemění viditelné pořadí.
\item Blok ČČČ a dvě nerozlišitelné modré kostky tvoří tři objekty se dvěma stejnými: $3!/2!=3$.

\end{enumerate}

\end{zadaniReseniBox}

\poznamka{Orientační práce v hodině: 0--12 min problém MAMA; 12--24 min odvození; 24--40 min úlohy; 40--45 min kontrola.}

\poznamka{Pozor na odlišný význam opakování oproti variacím s opakováním. Žáci nesmějí zaměňovat pevně určené četnosti a libovolné opakované volby.}

\end{document}
